\documentclass[../../main.tex]{subfiles}
\begin{document}

\paragraph{【新課程】}
座標の定められた空間において，
直線$l$は2点$(1,1,0)$，$(2,1,1)$を通り，
直線$m$は2点$(1,1,1)$，$(1,3,2)$を通る．
\begin{enumerate}[(1)]
  \item $l$を含み$m$に平行な平面の方程式を$ax+by+cz+d=0$の形に表せ．
  \item 点$(2,0,1)$を通り$l$，$m$の両方と交わる直線を$n$とする．
    $l$と$n$の交点および$m$と$n$の交点を求めよ．
\end{enumerate}

\paragraph{【旧課程】}
$a$を正の定数とし，$0 \leqq x \leqq 2\pi$の範囲で
関数$f(x)=-\cos x-\dfrac{a}{4}\cos 2x$を考える．
\begin{enumerate}[(1)]
  \item $f(x)$の最大値，最小値を求めよ．
  \item $a$の値の変化に応じて，$f(x)$のグラフの変曲点の個数はどのように変化するか．
    また，この個数の変化に応じて，$f(x)$のグラフの凹凸はどのように変化するか．
\end{enumerate}

\end{document}
